% !TEX root = ../main.tex
\section{Ejercicios resueltos}
En esta sección implementamos, ejecutamos y auditamos los tres ejercicios resueltos de la sección 8 de la guía oficial mediante el script \path{python/ejercicios_resueltos.py}. Para cada caso presentamos su deducción matemática, el código desarrollado y las salidas obtenidas.

\subsection{Resuelto 1: Convolución directa de polinomios}
Multiplicamos los polinomios $A(x) = 1 + 2x + 3x^2$ y $B(x) = 4 + 5x$, representados por los vectores de coeficientes $\mathbf{a} = [1, 2, 3]$ y $\mathbf{b} = [4, 5]$ ($m=3, n=2$). La longitud lineal resultante es $m+n-1 = 4$. El cálculo directo de productos cruzados confirma los coeficientes:
\begin{align}
c_0 &= a_0 b_0 = 1 \cdot 4 = 4, \\
c_1 &= a_0 b_1 + a_1 b_0 = 1 \cdot 5 + 2 \cdot 4 = 13, \\
c_2 &= a_1 b_1 + a_2 b_0 = 2 \cdot 5 + 3 \cdot 4 = 22, \\
c_3 &= a_2 b_1 = 3 \cdot 5 = 15.
\end{align}
El vector producto resultante es $C = [4, 13, 22, 15]$. Código fuente y salida ejecutada:
\begin{lstlisting}[language=Python,caption={Ejercicio resuelto 1: convolución directa}]
def resuelto_1():
    a, b = [1, 2, 3], [4, 5]
    result = naive_convolution(a, b)
    assert result == [4, 13, 22, 15]
    return {"a": a, "b": b, "c": result}
\end{lstlisting}
\noindent\textbf{Resultado de la ejecución:}
\begin{lstlisting}[style=consolestyle]
>>> resuelto_1()
{'a': [1, 2, 3], 'b': [4, 5], 'c': [4, 13, 22, 15]}
\end{lstlisting}


\subsection{Resuelto 2: FFT, IFFT y multiplicación con zero padding}
Para multiplicar $A(x)$ y $B(x)$ mediante FFT, ajustamos la longitud a la potencia de dos inmediata $N = 2^{\lceil \log_2(4) \rceil} = 4$. Aplicamos zero padding a $\mathbf{a}$ obteniendo $\mathbf{a}_{\text{pad}} = [1, 2, 3, 0]$ y a $\mathbf{b}$ obteniendo $\mathbf{b}_{\text{pad}} = [4, 5, 0, 0]$.

La evaluación en las 4 raíces cuartas de la unidad $\{\omega_4^0 = 1, \omega_4^1 = i, \omega_4^2 = -1, \omega_4^3 = -i\}$ produce el espectro:
\begin{align}
A_0 &= 1 + 2(1) + 3(1)^2 + 0 = 6, \\
A_1 &= 1 + 2(i) + 3(i)^2 + 0 = -2 + 2i, \\
A_2 &= 1 + 2(-1) + 3(-1)^2 + 0 = 2, \\
A_3 &= 1 + 2(-i) + 3(-i)^2 + 0 = -2 - 2i.
\end{align}
Comprobamos que la Transformada Inversa $\operatorname{IFFT}(\operatorname{FFT}(\mathbf{a}_{\text{pad}}))$ recupera el vector original con un residuo numérico menor a $10^{-16}$. El producto punto a punto en frecuencia seguido de IFFT recupera con exactitud los coeficientes $[4, 13, 22, 15]$, con un residuo imaginario residual de $6.66 \times 10^{-16}$, cumpliendo con creces la tolerancia establecida de $10^{-7}$.
\begin{lstlisting}[language=Python,caption={Ejercicio resuelto 2: FFT, IFFT y multiplicación}]
def resuelto_2():
    a, b = [1, 2, 3], [4, 5]
    padded = a + [0]
    roundtrip = ifft(fft(padded))
    result = multiply_fft(a, b)
    assert result == [4, 13, 22, 15]
    assert max(abs(x - y) for x, y in zip(padded, roundtrip)) < 1e-7
    return {"fft": fft(padded), "ifft_fft": roundtrip,
            "producto": result, "error_imaginario": max(abs(z.imag) for z in
             convolution_fft_raw(a, b))}
\end{lstlisting}
\noindent\textbf{Resultado de la ejecución:}
\begin{lstlisting}[style=consolestyle]
>>> resuelto_2()
{'fft': [(6+0j), (-1.9999999999999998+2j), (2+0j), (-2-2j)],
 'ifft_fft': [(1-0j),
              (2+5.721188726109833e-18j),
              (3-0j),
              -5.721188726109833e-18j],
 'producto': [4, 13, 22, 15],
 'error_imaginario': 6.661338147750939e-16}
\end{lstlisting}


\subsection{Resuelto 3: NTT e INTT directa sobre el campo finito $\mathbb{F}_{17}$}
Evaluamos el vector $\mathbf{a} = [1, 2, 3, 4]$ de longitud $N=4$ sobre el campo primo $\mathbb{F}_{17}$. Como $17 - 1 = 16$ y $4 \mid 16$, la raíz cuarta fijada en la guía es $\omega = 4$, la cual verifica:
\begin{equation}
4^1 \equiv 4, \quad 4^2 \equiv 16 \equiv -1, \quad 4^3 \equiv 64 \equiv 13, \quad 4^4 \equiv 256 \equiv 1 \pmod{17}.
\end{equation}
Como $4^4 \equiv 1$ y $4^2 \not\equiv 1 \pmod{17}$, $\omega = 4$ tiene orden multiplicativo exacto 4. La evaluación discreta directa:
\begin{equation}
A_k = \sum_{j=0}^3 a_j \cdot 4^{jk} \pmod{17},
\end{equation}
produce los valores espectrales:
\begin{align}
A_0 &= (1 + 2 + 3 + 4) \bmod 17 = 10 \bmod 17 = 10, \\
A_1 &= (1 + 2 \cdot 4 + 3 \cdot 16 + 4 \cdot 13) \bmod 17 = 109 \bmod 17 = 7, \\
A_2 &= (1 + 2 \cdot 16 + 3 \cdot 1 + 4 \cdot 16) \bmod 17 = 100 \bmod 17 = 15, \\
A_3 &= (1 + 2 \cdot 13 + 3 \cdot 16 + 4 \cdot 4) \bmod 17 = 91 \bmod 17 = 6.
\end{align}
Obtenemos el espectro NTT $[10, 7, 15, 6]$. Para la transformada inversa (INTT), calculamos la raíz inversa $\omega^{-1} \equiv 4^{-1} \equiv 13 \pmod{17}$ y el factor de escala $N^{-1} \equiv 4^{-1} \equiv 13 \pmod{17}$, recuperando el vector original $[1, 2, 3, 4]$ con exactitud aritmética absoluta.
\begin{lstlisting}[language=Python,caption={Ejercicio resuelto 3: NTT e INTT directa en F17}]
def resuelto_3():
    a = [1, 2, 3, 4]
    transformed = ntt_direct(a, 4, 17)
    recovered = intt_direct(transformed, 4, 17)
    assert transformed == [10, 7, 15, 6] and recovered == a
    return {"entrada": a, "raiz": 4, "modulo": 17,
            "ntt": transformed, "intt": recovered}
\end{lstlisting}
\noindent\textbf{Resultado de la ejecución:}
\begin{lstlisting}[style=consolestyle]
>>> resuelto_3()
{'entrada': [1, 2, 3, 4],
 'raiz': 4,
 'modulo': 17,
 'ntt': [10, 7, 15, 6],
 'intt': [1, 2, 3, 4]}
\end{lstlisting}

